\section{The Case-Study App: Accessible Podcast}
\label{sec:casestudy}

\subsection{Overview}

We built \textbf{Accessible Podcast}, a complete front-end Compose app, to ground the techniques. An
in-memory mock implements the \code{PodcastRepository} interface, so the UI code is identical to a real-backend version but runs anywhere with no configuration.

\subsection{User flow}

The app implements the classic three-step listening flow, plus a resume shortcut
(Figure~\ref{fig:flow}).

\begin{figure}[h]
	\centering
	\begin{tikzpicture}[
		s/.style={rounded corners=6pt,draw=brandblue,line width=1pt,fill=softbg,
				minimum width=3.1cm,minimum height=1.1cm,align=center,font=\small},
		a/.style={-{Stealth[length=2.5mm]},draw=brandblue,line width=1pt,font=\scriptsize}
		]
		\node[s] (h) {\textbf{Discover}\\{\scriptsize list of shows}};
		\node[s,right=2.2cm of h] (d) {\textbf{Podcast detail}\\{\scriptsize episodes}};
		\node[s,right=2.2cm of d] (p) {\textbf{Now Playing}\\{\scriptsize transport controls}};
		\draw[a] (h) -- node[above]{tap card} (d);
		\draw[a] (d) -- node[above]{tap episode} (p);
		\draw[a] (h.south) to[bend right=28] node[below,pos=0.5]{``Continue listening''} (p.south);
	\end{tikzpicture}
	\caption{User flow. Each transition is reachable by TalkBack double-tap and by Switch Access.}
	\label{fig:flow}
\end{figure}

\subsection{Architecture}

The code is organised into a thin data layer and a Compose UI layer
(Figure~\ref{fig:arch}). The UI depends
only on the \code{PodcastRepository} interface, never on the mock, thus we can substitute a real backend in without changing the UI's declaration and logic.

\begin{figure}[h]
	\centering
	\begin{tikzpicture}[
		layer/.style={rounded corners=5pt,draw,line width=0.8pt,minimum height=1cm,align=center,font=\small,text width=3.4cm},
		ui/.style={layer,draw=brandblue,fill=brandblue!7},
		dm/.style={layer,draw=brandgreen,fill=brandgreen!7},
		e/.style={-{Stealth[length=2mm]},draw=slate}
		]
		\node[ui] (nav) at (0,0) {\textbf{navigation}\\{\scriptsize NavHost: Home/Detail/Player}};
		\node[ui] (home) at (-3.6,-1.7) {\textbf{home}\\{\scriptsize HomeScreen, PodcastCard}};
		\node[ui] (det) at (0,-1.7) {\textbf{detail}\\{\scriptsize DetailScreen, EpisodeRow}};
		\node[ui] (play) at (3.6,-1.7) {\textbf{player}\\{\scriptsize PlayerScreen, PlayToggle}};
		\node[dm] (repo) at (0,-3.6) {\textbf{PodcastRepository}\\{\scriptsize interface}};
		\node[dm] (mock) at (-3.6,-3.6) {\textbf{MockPodcast\-Repository}};
		\node[dm] (model) at (3.6,-3.6) {\textbf{model}\\{\scriptsize Podcast, Episode}};
		\draw[e] (nav)--(home); \draw[e] (nav)--(det); \draw[e] (nav)--(play);
		\draw[e] (home)--(repo); \draw[e] (det)--(repo); \draw[e] (play)--(repo);
		\draw[e,dashed,brandgreen] (mock)-- node[below,font=\scriptsize]{implements} (repo);
		\draw[e] (repo)--(model);
	\end{tikzpicture}
	\caption{Package architecture. The Compose UI layer (blue) talks to a single repository interface;
		the mock (green) is the only implementation and can be swapped for a real one with no UI change.}
	\label{fig:arch}
\end{figure}

\subsection{The mock data layer}

The repository exposes an observable catalogue as a \code{StateFlow}. Toggling a bookmark mutates the
flow, so every screen recomposes exactly as it would against a reactive backend, which also lets
us verify that state changes are announced to TalkBack.

\begin{lstlisting}[caption={The repository interface the whole UI depends on.},captionpos=b]
interface PodcastRepository {
    val podcasts: StateFlow<List<Podcast>>   // observable catalogue
    fun getEpisode(id: String): Episode?
    fun toggleBookmark(episodeId: String)     // mutates state -> recomposition
    fun toggleDownload(episodeId: String)
}
\end{lstlisting}

\subsection{Where each technique lives}

Table~\ref{tab:map} maps every accessibility technique discussed in this report to the exact file in
the project, so the demo and the prose stay in lock-step.

\begin{table}[h]
	\centering
	\small
	\begin{tabularx}{\textwidth}{@{}l l Y@{}}
		\toprule
		\textbf{Report section}   & \textbf{Technique}                                 & \textbf{File}                                   \\
		\midrule
		Core (§\ref{sec:core})    & \code{contentDescription}, decorative \code{null}  & \texttt{ui/components/CoverArt.kt}              \\
		Core                      & \code{Role.Button} + \code{stateDescription}       & \texttt{ui/player/PlayToggle.kt}                \\
		Core                      & 48\,dp touch targets                               & \texttt{ui/home/PodcastCard.kt}                 \\
		Core                      & \code{sp} typography, font scaling                 & \texttt{ui/theme/Type.kt}                       \\
		Core                      & heading navigation                                 & \texttt{ui/home/HomeScreen.kt}                  \\
		Advanced (§\ref{sec:adv}) & \code{mergeDescendants}                            & \texttt{PodcastCard.kt}, \texttt{EpisodeRow.kt} \\
		Advanced                  & \code{clearAndSetSemantics} + \code{customActions} & \texttt{EpisodeRow.kt}                          \\
		Advanced                  & \code{hideFromAccessibility}                       & \texttt{CoverArt.kt}                            \\
		Advanced                  & \code{isTraversalGroup} + \code{traversalIndex}    & \texttt{HomeScreen.kt}                          \\
		Background                & \code{liveRegion} announcements                    & \texttt{ui/player/PlayerScreen.kt}              \\
		Testing (§\ref{sec:test}) & Compose assertions, ATF                            & \texttt{androidTest/AccessibilityTest.kt}       \\
		\bottomrule
	\end{tabularx}
	\caption{Technique-to-file map. Every claim in the report is backed by runnable code.}
	\label{tab:map}
\end{table}
